\documentclass[10pt,a4paper]{amsart}
\usepackage[margin=19mm]{geometry}
\usepackage{amsmath,amssymb,mathtools}
\usepackage[scr=rsfs,cal=euler]{mathalfa}
\usepackage{xfrac}
\usepackage[unicode,hidelinks,bookmarks=false]{hyperref}
\newcommand{\ie}{i.e.\ }
\newcommand{\cf}{cf.\ }
\newcommand{\resp}{respectively\ }
\newcommand{\Subsection}{\subsection}
\providecommand{\nobreakdash}{\nobreak\hbox{-}}
\setlength{\emergencystretch}{2em}
\multlinegap0pt
\usepackage{amssymb}
\usepackage{xcolor}
\newtheorem{proposition}{Proposition}
\newcommand{\Sp}{\mathfrak{S}}
\newcommand{\T}{\mathbb{T}}
\title{Global solutions for the Alber equation \\ in $H^1\Sp^1(\mathbb{T})$}
\author{Agissilaos Athanassoulis}
\date{Source: arXiv:2604.16998v2 (CC BY 4.0)}
\begin{document}
\maketitle

\noindent\emph{Source and licence.} Global solutions for the Alber equation \\ in $H^1\Sp^1(\mathbb{T})$. Version arXiv:2604.16998v2 (CC BY 4.0), distributed by arXiv under the Creative Commons Attribution 4.0 International licence, http://creativecommons.org/licenses/by/4.0/, as recorded in the arXiv licence metadata for this exact version and shown on the abstract page https://arxiv.org/abs/2604.16998v2. Full licence notice: \texttt{licenses/arxiv-2604.16998-CC-BY-4.0.txt}.\par\smallskip

\noindent\textit{Editorial scope.} Counted unit: the article's proposition prop:local (source0.tex lines 752--761). The article's complete original proof of it is delivered with it (source0.tex lines 763--829, one \texttt{proof} environment of 67 lines). No mathematics is written, repaired or re-derived: the statement and the proof are the article's own lines, in the article's own notation, and no retained line is substituted or re-flowed.  Retained uncounted as the source-local context the proof needs: lines 435--444; lines 662--675; lines 746--750.  Nothing was omitted from any retained span, so the package carries no omission record.  No retained line contains a citation, so the wrapper carries no bibliography and no bibliography entry.  No retained line contains a figure environment, an image-inclusion command or drawing code, so no image is delivered and the package declares no asset dependency.  Editorial formatting only, in addition: an AMS \texttt{amsart} preamble that restates the article's own declarations and macros used by the retained lines, namely the environments \texttt{proposition} on separate counters, together with the standard packages those lines need.  The retained environments print in this document as Proposition 1 in order of appearance, whereas the article prints the same items with the numbering of its own full document.  The complete native source of the article as distributed in the arXiv e-print for this exact version is bundled unchanged as \texttt{sources/198711/source0.tex}.
\begin{align}
  \|S(t)\gamma_0\|_{L^\infty([0,T];\,H^s\Sp^r)}
  &= \|\gamma_0\|_{H^s\Sp^r},
  \label{eq:free_iso}\\[4pt]
  \biggl\|\int_0^t S(t-\tau)F(\tau)\,d\tau
  \biggr\|_{L^\infty([0,T];\,H^s\Sp^1)}
  &\le T\,
  \|F\|_{L^\infty([0,T];\,H^s\Sp^1)}.
  \label{eq:Duh_Linfty}
\end{align}

\begin{proposition}[Bilinear bound]
\label{prop:bilinear}
Let $s>1/2$, $d=1$. For $\gamma_1,\gamma_2\in H^s\Sp^1(\T)$,
\begin{equation}\label{eq:bilinear}
  \|[V_{\rho_{\gamma_1}},\gamma_2]\|_{H^s\Sp^1}
  \le C\,\|\gamma_1\|_{H^s\Sp^1}
  \,\|\gamma_2\|_{H^s\Sp^1}.
\end{equation}
In particular, the nonlinearity $\mathcal{N}(\gamma)=[V_{\rho_\gamma},\gamma]$ satisfies
\begin{equation}\label{eq:nonlin_bound}
  \|\mathcal{N}(\gamma)\|_{H^s\Sp^1}
  \le C\,\|\gamma\|_{H^s\Sp^1}^2.
\end{equation}
\end{proposition}

\begin{equation}\label{eq:Duhamel_map}
  \Phi(\gamma)(t) \;:=\; S(t)\gamma_0
  - iq\int_0^t S(t-\tau)\,\mathcal{N}(\gamma(\tau))\,d\tau,
  \qquad t\in[0,T].
\end{equation}

\begin{proposition}[Local well-posedness on $H^1\Sp^1$]
\label{prop:local}
For every $\delta>0$ there exists $T>0$ such that if $\|\gamma_0\|_{H^1\Sp^1}\le\delta$, then the Duhamel map $\Phi_{\gamma_0}$ is a contraction on $\bigl\{\gamma\in X_T : \|\gamma\|_{X_T}\le 2\delta\bigr\}$.

Moreover, the solution depends Lipschitz-continuously on the initial datum: if $\|\gamma_0^{(1)}\|_{H^1\Sp^1}, \|\gamma_0^{(2)}\|_{H^1\Sp^1}\le\delta$ and $\gamma^{(1)},\gamma^{(2)}$ are the corresponding solutions, then
\begin{equation}\label{eq:lwp_stab}
  \|\gamma^{(1)}-\gamma^{(2)}\|_{X_T}
  \le 2\,\|\gamma_0^{(1)}-\gamma_0^{(2)}\|_{H^1\Sp^1}.
\end{equation}
\end{proposition}

\begin{proof}
Write $\delta=\|\gamma_0\|_{H^1\Sp^1}$.

\medskip

\noindent\textbf{Step 1: Linear estimate.} By isometry of the free evolution,
\begin{equation}\label{eq:linear_XT}
  \|S(t)\gamma_0\|_{X_T}
  = \|\gamma_0\|_{H^1\Sp^1}
  = \delta.
\end{equation}

\medskip

\noindent\textbf{Step 2: Nonlinear estimate.} By Proposition~\ref{prop:bilinear} and the Duhamel estimate~\eqref{eq:Duh_Linfty}:
\begin{equation}\label{eq:Duh_S1}
  \biggl\|\int_0^t S(t-\tau)
  \mathcal{N}(\gamma)\,d\tau
  \biggr\|_{X_T}
  \le C\,T\,\|\gamma\|_{X_T}^2.
\end{equation}

\noindent\textbf{Step 3: Closing.} On the ball $\|\gamma\|_{X_T}\le 2\delta$:
\[
  \|\Phi(\gamma)\|_{X_T}
  \le \delta + C\,T\,(2\delta)^2
  = \delta\,(1 + 4C\,T\,\delta).
\]
Choosing $T$ small enough that $4C\,T\,\delta\le 1/2$, i.e.\
\begin{equation}\label{eq:delta_T}
  T \le \frac{1}{8C\,\delta},
\end{equation}
gives $\|\Phi(\gamma)\|_{X_T}\le\tfrac{3}{2}\delta<2\delta$, so $\Phi$ maps the ball $\|\gamma\|_{X_T}\le 2\delta$ to itself.

For the contraction, let $\gamma_1,\gamma_2$ lie in the ball. Using the linearity $\rho_{\gamma_1}-\rho_{\gamma_2}=\rho_{\gamma_1-\gamma_2}$, polarize:
\[
  \mathcal{N}(\gamma_1)-\mathcal{N}(\gamma_2)
  = [V_{\rho_{\gamma_1-\gamma_2}},\gamma_1]
  + [V_{\rho_{\gamma_2}},\gamma_1-\gamma_2].
\]
Applying Proposition~\ref{prop:bilinear} to each term gives
\[
  \|\mathcal{N}(\gamma_1)-\mathcal{N}(\gamma_2)\|_{H^1\Sp^1}
  \le C\bigl(\|\gamma_1\|_{H^1\Sp^1}+\|\gamma_2\|_{H^1\Sp^1}\bigr)
  \,\|\gamma_1-\gamma_2\|_{H^1\Sp^1},
\]
and the same Duhamel estimate as in Step~2 yields
\[
  \|\Phi(\gamma_1)-\Phi(\gamma_2)\|_{X_T}
  \le 4C\,T\,\delta\,
  \|\gamma_1-\gamma_2\|_{X_T}
  \le \tfrac{1}{2}\,
  \|\gamma_1-\gamma_2\|_{X_T},
\]
so $\Phi$ is a contraction.

\medskip

\noindent\textbf{Step 4: Continuous dependence.} Let $\gamma^{(j)}=\Phi_{\gamma_0^{(j)}}(\gamma^{(j)})$ for $j=1,2,$ with $\Phi_{\gamma_0}$ as in~\eqref{eq:Duhamel_map}. Subtracting and using the linear estimate together with the Lipschitz bound on $\mathcal{N}$ from Step~3,
\[
  \|\gamma^{(1)}-\gamma^{(2)}\|_{X_T}
  \le \|\gamma_0^{(1)}-\gamma_0^{(2)}\|_{H^1\Sp^1}
  + \tfrac12
    \|\gamma^{(1)}-\gamma^{(2)}\|_{X_T},
\]
yielding~\eqref{eq:lwp_stab}.
\end{proof}

\end{document}
